\documentclass[11pt,a4paper]{article}
% Wspólny szablon dla kart supermocy (jedna karta = jedna strona A4)
% Użycie w karcie: \documentclass[11pt,a4paper]{article} % Wspólny szablon dla kart supermocy (jedna karta = jedna strona A4)
% Użycie w karcie: \documentclass[11pt,a4paper]{article} % Wspólny szablon dla kart supermocy (jedna karta = jedna strona A4)
% Użycie w karcie: \documentclass[11pt,a4paper]{article} \input{../preambula} \begin{document}

% Silnik tectonic (XeTeX): polskie znaki przez fontspec, nie przez inputenc/T1
\usepackage{fontspec}
\setmainfont{Arial}
\usepackage[polish]{babel}
\usepackage[a4paper,margin=1.5cm]{geometry}
\usepackage{amsmath,amssymb}
\usepackage{tikz}
\usetikzlibrary{arrows.meta,positioning,calc,shapes.geometric,decorations.pathreplacing}
\usepackage[most]{tcolorbox}
\usepackage{enumitem}
\usepackage{booktabs}
\usepackage{xcolor}

\pagestyle{empty}
\setlength{\parindent}{0pt}
\setlength{\parskip}{0.4em}

\definecolor{cel}{HTML}{1F4E79}
\definecolor{regula}{HTML}{E8F1FA}
\definecolor{pulapka}{HTML}{FDECEA}
\definecolor{okres}{HTML}{C0392B}
\definecolor{skonczony}{HTML}{1E8449}

% Ramka na regułę
\newtcolorbox{regula}[1][]{
  colback=regula, colframe=cel, boxrule=0.8pt, arc=2pt,
  left=6pt, right=6pt, top=4pt, bottom=4pt, title=#1,
  fonttitle=\bfseries, coltitle=white
}

% Ramka na pułapkę z kartkówki
\newtcolorbox{pulapka}[1][]{
  colback=pulapka, colframe=okres, boxrule=0.8pt, arc=2pt,
  left=6pt, right=6pt, top=4pt, bottom=4pt, title=#1,
  fonttitle=\bfseries, coltitle=white
}

% Nagłówek karty: numer, tytuł, cel
\newcommand{\naglowekKarty}[3]{%
  {\LARGE\bfseries\color{cel} Karta #1: #2}\\[2pt]
  {\large Cel: #3}\par
  \vspace{0.3em}\noindent\rule{\linewidth}{1pt}\par
}

% Pole na obliczenia ucznia: ramka z delikatną kratką, argument = wysokość
\newcommand{\miejsce}[1]{%
  \par\vspace{0.2em}\noindent
  \begin{tikzpicture}
    \draw[step=0.5cm, gray!25, very thin] (0,0) grid (\textwidth,#1);
    \draw[gray!60, rounded corners=2pt] (0,0) rectangle (\textwidth,#1);
    \node[anchor=north west, font=\scriptsize, text=gray] at (0.05,#1) {obliczenia};
  \end{tikzpicture}\par}
 \begin{document}

% Silnik tectonic (XeTeX): polskie znaki przez fontspec, nie przez inputenc/T1
\usepackage{fontspec}
\setmainfont{Arial}
\usepackage[polish]{babel}
\usepackage[a4paper,margin=1.5cm]{geometry}
\usepackage{amsmath,amssymb}
\usepackage{tikz}
\usetikzlibrary{arrows.meta,positioning,calc,shapes.geometric,decorations.pathreplacing}
\usepackage[most]{tcolorbox}
\usepackage{enumitem}
\usepackage{booktabs}
\usepackage{xcolor}

\pagestyle{empty}
\setlength{\parindent}{0pt}
\setlength{\parskip}{0.4em}

\definecolor{cel}{HTML}{1F4E79}
\definecolor{regula}{HTML}{E8F1FA}
\definecolor{pulapka}{HTML}{FDECEA}
\definecolor{okres}{HTML}{C0392B}
\definecolor{skonczony}{HTML}{1E8449}

% Ramka na regułę
\newtcolorbox{regula}[1][]{
  colback=regula, colframe=cel, boxrule=0.8pt, arc=2pt,
  left=6pt, right=6pt, top=4pt, bottom=4pt, title=#1,
  fonttitle=\bfseries, coltitle=white
}

% Ramka na pułapkę z kartkówki
\newtcolorbox{pulapka}[1][]{
  colback=pulapka, colframe=okres, boxrule=0.8pt, arc=2pt,
  left=6pt, right=6pt, top=4pt, bottom=4pt, title=#1,
  fonttitle=\bfseries, coltitle=white
}

% Nagłówek karty: numer, tytuł, cel
\newcommand{\naglowekKarty}[3]{%
  {\LARGE\bfseries\color{cel} Karta #1: #2}\\[2pt]
  {\large Cel: #3}\par
  \vspace{0.3em}\noindent\rule{\linewidth}{1pt}\par
}

% Pole na obliczenia ucznia: ramka z delikatną kratką, argument = wysokość
\newcommand{\miejsce}[1]{%
  \par\vspace{0.2em}\noindent
  \begin{tikzpicture}
    \draw[step=0.5cm, gray!25, very thin] (0,0) grid (\textwidth,#1);
    \draw[gray!60, rounded corners=2pt] (0,0) rectangle (\textwidth,#1);
    \node[anchor=north west, font=\scriptsize, text=gray] at (0.05,#1) {obliczenia};
  \end{tikzpicture}\par}
 \begin{document}

% Silnik tectonic (XeTeX): polskie znaki przez fontspec, nie przez inputenc/T1
\usepackage{fontspec}
\setmainfont{Arial}
\usepackage[polish]{babel}
\usepackage[a4paper,margin=1.5cm]{geometry}
\usepackage{amsmath,amssymb}
\usepackage{tikz}
\usetikzlibrary{arrows.meta,positioning,calc,shapes.geometric,decorations.pathreplacing}
\usepackage[most]{tcolorbox}
\usepackage{enumitem}
\usepackage{booktabs}
\usepackage{xcolor}

\pagestyle{empty}
\setlength{\parindent}{0pt}
\setlength{\parskip}{0.4em}

\definecolor{cel}{HTML}{1F4E79}
\definecolor{regula}{HTML}{E8F1FA}
\definecolor{pulapka}{HTML}{FDECEA}
\definecolor{okres}{HTML}{C0392B}
\definecolor{skonczony}{HTML}{1E8449}

% Ramka na regułę
\newtcolorbox{regula}[1][]{
  colback=regula, colframe=cel, boxrule=0.8pt, arc=2pt,
  left=6pt, right=6pt, top=4pt, bottom=4pt, title=#1,
  fonttitle=\bfseries, coltitle=white
}

% Ramka na pułapkę z kartkówki
\newtcolorbox{pulapka}[1][]{
  colback=pulapka, colframe=okres, boxrule=0.8pt, arc=2pt,
  left=6pt, right=6pt, top=4pt, bottom=4pt, title=#1,
  fonttitle=\bfseries, coltitle=white
}

% Nagłówek karty: numer, tytuł, cel
\newcommand{\naglowekKarty}[3]{%
  {\LARGE\bfseries\color{cel} Karta #1: #2}\\[2pt]
  {\large Cel: #3}\par
  \vspace{0.3em}\noindent\rule{\linewidth}{1pt}\par
}

% Pole na obliczenia ucznia: ramka z delikatną kratką, argument = wysokość
\newcommand{\miejsce}[1]{%
  \par\vspace{0.2em}\noindent
  \begin{tikzpicture}
    \draw[step=0.5cm, gray!25, very thin] (0,0) grid (\textwidth,#1);
    \draw[gray!60, rounded corners=2pt] (0,0) rectangle (\textwidth,#1);
    \node[anchor=north west, font=\scriptsize, text=gray] at (0.05,#1) {obliczenia};
  \end{tikzpicture}\par}


\begin{document}

\naglowekKarty{4}{Sprawdź się: punkty procentowe}{5 minut, pp czy procenty}

\vspace{0.2em}
Imię: Nadia \hfill Data: \rule{3cm}{0.4pt} \hfill Punkty: \rule{1.2cm}{0.4pt} / 6 pkt

{\small W każdym zadaniu pokaż obliczenia. Sama odpowiedź bez obliczeń nie daje pełnej liczby punktów.}

\vspace{0.4em}

\textbf{Zadanie 1} (2 pkt). Frekwencja w klasie wzrosła z $60\%$ do $75\%$. O ile punktów procentowych wzrosła? Pokaż obliczenia.
\miejsce{2.5cm}
Odpowiedź: $\rule{3cm}{0.4pt}$ pp

\vspace{0.4em}
\textbf{Zadanie 2} (2 pkt). Udział uczniów z ocenami celującymi wzrósł z $20\%$ do $30\%$. Pokaż obliczenia.
\begin{enumerate}[label=\alph*), leftmargin=*, itemsep=2pt]
  \item O ile punktów procentowych? \rule{2.5cm}{0.4pt} \hspace{0.5cm} b) O ile procent? \rule{2.5cm}{0.4pt}
\end{enumerate}
\miejsce{3.5cm}

\textbf{Zadanie 3} (2 pkt). Marek twierdzi: „Poparcie wzrosło z $20\%$ do $25\%$, czyli o $5\%$”. Wyjaśnij jego błąd, a potem oblicz, o ile procent naprawdę wzrosło poparcie. Zapisz tok rozumowania.
\miejsce{4.5cm}
Odpowiedź: \rule{4cm}{0.4pt}

\vfill
\begin{tcolorbox}[colback=white, colframe=cel, boxrule=0.6pt, arc=2pt]
Ranga: \quad $\square$ 5--6 pkt: \textbf{Mistrzyni procentów!} \quad $\square$ 3--4 pkt: \textbf{Bystre oko!} \quad $\square$ 1--2 pkt: \textbf{Krok do przodu!}
\end{tcolorbox}

\end{document}
